\documentclass[10pt]{article}
\usepackage[T1]{fontenc}
\usepackage[utf8]{inputenc}
\usepackage[french]{babel}
\usepackage{amsmath,amssymb,mathrsfs,mathpazo}
\usepackage{geometry}
\geometry{a4paper,margin=2cm}
\usepackage{graphicx,array,tabularx,tabularray,colortbl,multirow,multicol,booktabs,enumitem,pifont,ulem,fancybox,diagbox,cancel}
\usepackage{pstricks,pstricks-add,pst-plot,pst-tree,pst-node,pst-text,pst-func,pst-math,pst-bspline,pst-3dplot,pst-eucl,pst-solides3d}
\let\pscalculate\undefined
\usepackage{pst-bezier}
\usepackage{tikz,pgfplots,tkz-tab}
\usetikzlibrary{arrows,arrows.meta,patterns,calc,positioning,trees,decorations.markings,shapes.geometric,shapes.misc,shapes.arrows,intersections,decorations.pathmorphing,decorations.pathreplacing,backgrounds,shadows,babel}
\pgfplotsset{compat=1.18}
\usepgfplotslibrary{fillbetween}
\usepackage[np]{numprint}
\usepackage{listings,makecell,pas-tableur,twoopt}
\let\barre\undefined
\usepackage[pstricks]{tablvar}
\definecolor{darkgreen}{RGB}{0,100,0}
\providecommand{\red}{\color{red}}
\providecommand{\blue}{\color{blue}}
\providecommand{\green}{\color{green}}
\providecommand{\black}{\color{black}}
\providecommand{\white}{\color{white}}
\providecommand{\gray}{\color{gray}}
\providecommand{\cyan}{\color{cyan}}
\providecommand{\magenta}{\color{magenta}}
\providecommand{\yellow}{\color{yellow}}
\providecommand{\violet}{\color{violet}}
\providecommand{\orange}{\color{orange}}
\providecommand{\euro}{\texteuro}
\providecommand{\EUR}{\texteuro}
\providecommand{\ieme}{\textsuperscript{e}}
\providecommand{\degres}{\ensuremath{^\circ}}
\providecommand{\degree}{\ensuremath{^\circ}}
\providecommand{\vv}[1]{\overrightarrow{#1}}
\providecommand{\Coord}[1]{#1}
\providecommand{\decosix}{}
\providecommand{\starredbullet}{\ensuremath{\bullet}}
\pagestyle{empty}
\setlength{\parindent}{0pt}
\usepackage{adjustbox}\DeclareUnicodeCharacter{2236}{:}\DeclareUnicodeCharacter{2190}{\ensuremath{\leftarrow}}
\providecommand{\R}{\mathbb{R}}
\providecommand{\N}{\mathbb{N}}
\providecommand{\Z}{\mathbb{Z}}
\providecommand{\Q}{\mathbb{Q}}
\providecommand{\D}{\mathbb{D}}
\providecommand{\C}{\mathbb{C}}
\providecommand{\e}{{\mathrm{e}}}
\providecommand{\vect}[1]{\overrightarrow{#1}}
\providecommand{\vectt}[1]{\overrightarrow{\mathrm{#1}}}
\providecommand{\vvt}[1]{\overrightarrow{\mathrm{#1}}}
\providecommand{\np}[2][]{{#2~#1}}
\providecommand{\vv}[1]{{\overrightarrow{#1}}}
\renewcommand{\d}{{\mathrm{d}}}
\providecommand{\degres}{{^\circ}}
\providecommand{\degre}{{^\circ}}
\providecommand{\blue}{\color{blue}}
\providecommand{\red}{\color{red}}
\providecommand{\enskip}{\quad}
\providecommand{\llbracket}{\lbrack\!\lbrack}
\providecommand{\rrbracket}{\rbrack\!\rbrack}
\providecommand{\euro}{{\text{€}}}
\providecommand{\no}{{\text{n°}}}
\providecommand{\interval}[2]{{[#1;#2]}}
\providecommand{\intervFO}[2]{{[#1;#2[}}
\providecommand{\dotfill}{\ldots}

\providecommand{\ds}{\displaystyle}
\providecommand{\textding}[1]{#1}
\providecommand{\Oij}{\(\mathrm{O};\vec{\imath},\vec{\jmath}\)}
\providecommand{\Oijk}{\(\mathrm{O};\vec{\imath},\vec{\jmath},\vec{k}\)}
\providecommand{\Ouv}{\(\mathrm{O};\vec{u},\vec{v}\)}
\providecommand{\intervFF}[2]{\ensuremath{[#1;#2]}}
\renewcommand{\euro}{\ifmmode\text{\texteuro}\else\texteuro\fi}
\DeclareUnicodeCharacter{00B0}{\ifmmode^{\circ}\else\textdegree\fi}
\DeclareUnicodeCharacter{20AC}{\ifmmode\text{\texteuro}\else\texteuro\fi}
% Énoncés : Éducation nationale. Transcription et corrigés : APMEP, Denis Vergès et contributeurs. Découpage et classement : CoopMaths, Sébastien Lozano, Éric Elter et contributeurs. Adaptation de la présentation : Maths Explorer.
% CC BY-SA : https://coopmaths.fr/www/about/
% https://coopmaths.fr/alea/static/bac/2021/tex/bac_2021_03_sujet1_4_cor.tex
\begin{document}
\noindent\textit{APMEP / CoopMaths. CC BY-SA. Présentation adaptée par Maths Explorer.}\par\medskip
\medskip

\begin{tabularx}{\linewidth}{|X|}\hline
\textbf{Principaux domaines abordés : Suites numériques; raisonnement par récurrence ; suites géométriques.}\\ \hline
\end{tabularx}

\medskip

La suite $\left(u_n\right)$ est définie sur $\N$ par $u_0 = 1$ et pour tout $n$, 
$u_{n+1} = \dfrac{3}{4}u_n + \dfrac{1}{4}n + 1$.




\begin{enumerate}
\item 
Pour $n=0$, $u_1=u_{0+1}=\dfrac{3}{4}u_0 + \dfrac{1}{4}\times 0 +1 =\dfrac{3}{4}\times 1  +1 = \dfrac{7}{4}$.

Pour $n=1$, $u_1=u_{1+1}=\dfrac{3}{4}u_1 + \dfrac{1}{4}\times 1+1 = \dfrac{3}{4}\times \dfrac{7}{4}+\dfrac{1}{4} +1 = \dfrac{41}{16}$. 

\end{enumerate}

\medskip

\parbox{0.5\linewidth}{L'extrait, reproduit ci-contre, d'une feuille de calcul réalisée avec un tableur présente les valeurs des premiers termes de la suite $\left(u_n\right)$.} \hfill
\parbox{0.35\linewidth}{
\begin{tabularx}{\linewidth}{|c|*{2}{>{\centering \arraybackslash}X|}}\hline
\cellcolor{lightgray}{} 	& \cellcolor{lightgray}A&\cellcolor{lightgray}B \\ \hline
\cellcolor{lightgray}1	&$n$&$u_n$\\ \hline
\cellcolor{lightgray}2 	&0	&1\\ \hline
\cellcolor{lightgray}3 	&1	&1,75\\ \hline
\cellcolor{lightgray}4 	&2	&\np{2,5625}\\ \hline
\cellcolor{lightgray}5 	&3	&\np{3,421875}\\ \hline
\cellcolor{lightgray}6 	&4	&\np{4,31640625}\\ \hline
\end{tabularx}}

\medskip

\begin{enumerate}[resume]
\item
	\begin{enumerate}
		\item La formule, étirée ensuite vers le bas, que l'on peut écrire dans la cellule B3 de la feuille de calcul pour obtenir les termes successifs de $\left(u_n\right)$ dans la colonne B est:
		
\text{= 3/4 * B2 + 1/4 * A2 +1}.		
		
		\item La suite $\left(u_n\right)$ semble croissante.
	\end{enumerate}
\item
	\begin{enumerate}
		\item Soit $\mathcal{P}_n$ la propriété: $n \leqslant u_n \leqslant n + 1$.
		
\begin{list}{\textbullet}{}
\item \textbf{Initialisation}

Pour $n=0$, $u_0=1$ et $0 \leqslant 1 \leqslant 1$ donc $\mathcal{P}_0$ est vraie.

\item \textbf{Hérédité}

On suppose $\mathcal{P}_n$ vraie, c'est-à-dire: $n \leqslant u_n \leqslant n+1$ (hypothèse de récurrence).

$\begin{gathered}n \leqslant u_n \leqslant n+1 \iff \dfrac{3}{4}n \leqslant \dfrac{3}{4}u_n \leqslant \dfrac{3}{4}\left (n+1\right )\\[7pt]
 \phantom{n \leqslant u_n \leqslant n+1}
 \iff \dfrac{3}{4}n +\dfrac{1}{4} n \leqslant \dfrac{3}{4}u_n  +\dfrac{1}{4} n  \leqslant \dfrac{3}{4}\left (n+1\right )  +\dfrac{1}{4} n \\[7pt]
 \phantom{n \leqslant u_n \leqslant n+1}
 \iff n \leqslant \dfrac{3}{4}u_n  +\dfrac{1}{4} n  \leqslant  n+\dfrac{3}{4}\\[7pt]
 \phantom{n \leqslant u_n \leqslant n+1}
 \iff n+1 \leqslant \dfrac{3}{4}u_n  +\dfrac{1}{4} n +1  \leqslant  n+\dfrac{3}{4} +1
 \iff n+1 \leqslant u_{n+1} \leqslant n+\dfrac{7}{4}\end{gathered}$ 
 
 donc $n+1 \leqslant u_{n+1} \leqslant n+2$.

On a démontré que la propriété était vraie au rang $n+1$.

\item \textbf{Conclusion}

La propriété est vraie au rang 0, et elle est héréditaire pour tout $n\geqslant 0$; d'après le principe de récurrence, la propriété est vraie pour tout $n\geqslant 0$.
\end{list}		

On a donc démontré que, pour tout entier naturel $n$, on a: $n \leqslant u_n \leqslant n + 1$.
		
		\item
\begin{list}{\textbullet}{D'après la question précédente:}
\item Pour tout $n$, $n\leqslant u_n \leqslant n+1$ donc $n+1 \leqslant u_{n+1} \leqslant n+2$ donc 

$n \leqslant u_n \leqslant n+1 \leqslant u_{n+1} \leqslant n+2$ d'où on tire $u_n \leqslant u_{n+1}$ ce qui démontre que la suite $(u_n)$ est croissante.
\item Pour tout $n$, $n\leqslant u_n$; or $\ds\lim_{n\to +\infty} n = +\infty$ donc, par comparaison, $\ds\lim_{n\to +\infty} u_n = +\infty$.
\end{list}		
		
		
		\item
Pour tout $n$, $n\leqslant u_n \leqslant n+1$ donc pour tout $n>0$, on a:
$1 \leqslant \dfrac{u_n}{n} \leqslant \dfrac{n+1}{n}$ c'est-à-dire: 
$1 \leqslant \dfrac{u_n}{n} \leqslant 1+\dfrac{1}{n}$.

$\ds\lim_{n \to +\infty} \dfrac{1}{n}=0$ donc $\ds\lim_{n \to +\infty} 1+\dfrac{1}{n}=1$

Donc, d'après le théorème des gendarmes: $\ds\lim_{n\to +\infty} \dfrac{u_n}{n}=1$.

	\end{enumerate}
\item  On désigne par $\left(v_n\right)$ la suite définie sur $\N$ par $v_n = u_n - n$
	\begin{enumerate}
		\item
Pour tout $n$, $v_n=u_n-n$ donc $u_n=v_n+n$.

$v_{n+1}=u_{n+1} - (n+1) = 	\dfrac{3}{4}u_n + \dfrac{1}{4}n+1 -n-1
=\dfrac{3}{4}\left (v_n+n\right ) -\dfrac{3}{4}n
=\dfrac{3}{4}v_n +\dfrac{3}{4}n -\dfrac{3}{4}n
=\dfrac{3}{4}v_n$

$v_0=u_0-0=1$

Donc la suite $(v_n)$ est géométrique de raison $q=\dfrac{3}{4}$ et de premier terme $v_0=1$.
		
		
		\item 
On en déduit que, pour tout $n$, $v_n=v_0\times q^n = \left (\dfrac{3}{4}\right )^n$.

Comme $u_n=v_n+n$, on a $u_n=\left (\dfrac{3}{4}\right )^n + n$.		
		
	\end{enumerate}
\end{enumerate}

\bigskip

\end{document}
